\documentclass[a4paper, 11pt]{article}
%\documentclass{article}
\usepackage[utf8]{inputenc}
\usepackage{hyperref}
\usepackage{calc}
%\usepackage{fourier} %% changes everything to sans-serif!
\usepackage{amssymb}
\usepackage{amsmath}
\usepackage{amsthm}
\usepackage{amsfonts}
\usepackage{mathtools}
\usepackage{textcomp}
%\usepackage{tikz}
\usepackage{tikz-cd}
\usetikzlibrary{cd}
\usepackage{dirtytalk}
\usepackage{titlesec}
\usepackage{verbatim}
\usepackage{pdfpages}
\usepackage{caption}
%%\captionsetup[figure]{labelfont=footnotesize}
\usepackage{hyperref}
\usepackage{parskip}

%% neu
\usepackage{physics}
\usepackage{accents} %% for multiple subscripts
\usepackage{fancyhdr} %% for header and footer
\usepackage[a4paper, left=2.5cm, right=2.5cm, top=2cm]{geometry}
\usepackage{graphicx} % Required for inserting images
%\usepackage[ngerman]{babel}
\usepackage{csquotes}
\MakeOuterQuote{"}

\title{Spiele}
\author{Catharina Hock, Stefanie Hövermann}
\date{August  2026}

\newcommand{\answerline}[1][0.5cm]{%
    \noindent\rule{#1}{0.4pt}
}

\begin{document}
\maketitle 
Mehr Spiele hier. https://www.begabungslotse.de/methodenkoffer/gruppenzusammenhalt-und-aktivierung/alle-wechseln-den-platz


\section{Kennenlernspiele}

\subsection{Zipp - Zapp}

Alle sitzen im Stuhlkreis. Einer ist in der Mitte. Dieser zeigt auf einen Mitspieler uns sagt Zipp oder Zapp. Bei Zipp muss der Namen des linken Spielers, bei Zapp muss der Name des rechten Spielers genannt werden (kann auch mal getauscht werden). Wenn der Angesprochene diesen Namen nicht innerhalb von 3-4 Sekunden genannt hat, dann muss dieser in die Mitte. Sagt der Spieler in der Mitte Zipp-Zapp müssen sich alle einen neuen Platz suchen. Wer keinen findet muss in die Mitte.


\subsection{Alle, die…}

\begin{itemize}
    \item Dauer: ca. 10 Minuten
    \item Material: Stühle
    \item Spielbeschreibung: Alle Spieler*innen sitzen im Stuhlkreis. Eine freiwillige Person steht in die Mitte. Die Person in der Mitte sagt nun: ``Alle, die … (z. B. gerne Pizza essen).'' Anschließend müssen alle, die gerne Pizza essen, die Plätze tauschen. Die Person aus der Mitte versucht sich einen Platz zu sichern und sich auf den Stuhl zu setzen. Die Person, die übrig bleibt, ist als Nächstes an der Reihe.
    \item Variation : Die Spielleitung moderiert und gibt die Anweisungen. 
    Die Person in der Mitte hat die Aufgabe, sich auf den freien Stuhl zu setzen. Das kann die Person, die rechts von dem freien Stuhl sitzt, dadurch verhindern, indem sie schnell einen beliebigen Namen aus der Gruppe ruft. Die genannte Person muss nun möglichst schnell auf dem freien Stuhl Platz nehmen. Dadurch wird ein neuer Stuhl frei. Wenn es der Person in der Mitte gelingt sich zu setzen, bevor ein neuer Name genannt wird, muss die Person, die nicht schnell genug reagiert hat, selbst in die Mitte.
\end{itemize}

\subsection{2 Wahrheiten, 1 Lüge}

\begin{itemize}
    \item Dauer: ca. 10 bis 30 Minuten
    \item Material: ggf. Stühle
    \item Spielbeschreibung: Alle Teilnehmer*innen überlegen sich drei persönliche Dinge, die sie über sich erzählen möchten, z. B. besondere Erlebnisse, ungewöhnliche Interessen oder Hobbys, originelle Eigenschaften etc. Zwei Fakten sollen dabei wahr und ein Fakt frei erfunden sein. Der Reihe nach tragen die Teilnehmer*innen vor. Die anderen aus der Gruppe dürfen raten und abstimmen, was davon gelogen ist. Das Ziel des Spiels ist es, etwas Besonderes über die anderen zu erfahren. 
    \item Varianten: An einer Tafel wird notiert, welche Person am häufigsten die Lügen richtig erkannt hat.
\end{itemize}

\subsection{Ahh. . . Du bist es}
\begin{itemize}
    \item Alle stellen sich in einem Kreis auf. Der Spielleiter / die Spielleiterin nennt seinen bzw. ihren Namen und macht dazu eine auf ihn oder sie zutreffende Geste und einen Laut / ein Geräusch. Beispiel: Ich bin Anne. Um zu beschreiben, dass ich gerne koche, rühre ich in einem imaginären Kochtopf und sage `Mmh'. Vivian hingegen spielt gerne Basketball. Also sagt sie: `Ich bin Vivian.' Und deutet dann einen Basketballwurf mit einem `plopp' an.
    \item Alle Beteiligten antworten mit wohlig-euphorischem `Ahh … du bist [Wiederholung des Namens]' und mit Nachahmung der Geste.
    \item Nach jeder individuellen Namensnennung plus Geste wird die gesamte bisherige Reihe `nachgeahmt' mit Nennung des Namens und Nachahmung der Geste der bisherigen Teilnehmenden (vgl. `Ich packe meinen Koffer'). Gemeinsames Ziel: Eine Runde mit Namen und Gesten und dann rückwärts. Hilft beim Namen- und Kennenlernen. 
\end{itemize}

\section{Spiele für Zwischendurch}

\subsection{Atomspiel}
\begin{itemize}
    \item Dauer:ca. 10 Minuten
    \item Material: kein Material notwendig
    \item Spielbeschreibung: Alle Spieler*innen laufen zur Musik frei durch den Raum. Die Spielleitung hält die Musik an und ruft z. B. ``Atom 3''. Jetzt müssen sich alle Schüler*innen möglichst schnell zu Dreiergruppen (Kreis mit Handfassung) formieren. Wer keine Gruppe mehr findet, muss eine Zusatzaufgabe absolvieren (z. B. fünf Hampelmänner oder zehn Kniebeugen). 
    \item Varianten: Alle Spieler*innen laufen kreuz und quer durch den Raum. Dann sagt die Spielleitung: ``Alle, die schwarze Schnürsenkel haben!'' oder ``Alle, die blonde Haare haben.'' Daraufhin müssen sich die Teilnehmer*innen, bei denen die Eigenschaft zutrifft, zusammenfinden und mit einem Finger berühren. Sie bilden ein ``Atom''.
\end{itemize}

\subsection{Amöbe}

\begin{itemize}
    \item Zunächst werden die Regeln eines regulären Schnick-Schnack-Schnuck-Spiels erläutert. Es gibt folgende drei Optionen: Stein, Schere, Papier. Stein schlägt Schere, Schere schlägt Papier, Papier schlägt Stein.
    \item Nun wird die Variante ``Evolution'' anhand dieser (unwissenschaftlichen) Metaphern erklärt 
    \item Stufe 1 ``Amöbe'': Zu Beginn sind alle ``Amöben''. Als Amöbe bewegt man sich in der Hocke, mit an den Körper gelegten Armen und ruft dabei laut ``Amöbe, Amöbe, Amöbe''. Siegt eine Amöbe im Schnick-Schnack-Schnuck gegen die andere, entwickelt sie sich zum Schmetterling. Die Verlierer-Amöbe bleibt Amöbe.
    \item Stufe 2: ``Schmetterling'': Als Schmetterling flattert man durch den Raum. Trifft ein Schmetterling auf einen anderen Schmetterling kann ein Schnick-Schnack-Schnuck-Duell starten. Der Gewinner-Schmetterling wird zum Dino, der unterlegene Schmetterling geht zurück zum Amöben-Stadium.

    \item Gleiches gilt für die nun folgenden Stufen.

    \item Stufe 3 ``Dino'': Als Dino läuft man wie ein T-Rex brüllend durch die Gruppe.
    \item Stufe 4 ``Gorilla'': Ein Gorilla trommelt sich auf die Brust, während er sich auf der Suche nach einem weiteren Gorilla durch den Raum bewegt.
    \item Stufe 5 ``Mensch'': Sobald ein Spieler, eine Spielerin die Stufe Mensch erreicht hat, ist das Spiel beendet. 
\end{itemize}

\subsection{Kotzendes Känguru}
\begin{itemize}
    \item Die Gruppe steht im Kreis. Zunächst müssen ausgewählte verschiedene Figuren (s.u.) eingeführt, d.h. vorgemacht, werden. Eine Figur besteht immer aus einer Person und der rechts und links benachbarten Person.
    Während des Spiels steht immer eine Person in der Mitte, zeigt auf eine Person in der Runde und ruft eine Figur aus. Diese Person und die beiden benachbarten Personen stellen dann schnellstmöglich die vorgegebene Figur dar. Passiert ein Fehler oder dauert es sehr lange, wird der Platz mit der Person in der Mitte getauscht.
    Zu Beginn reichen drei Figuren. Wenn das Spiel wiederholt wird, können weitere Figuren hinzugenommen werden.
    \item Figuren 
    \item Kotzendes Känguru: Person in der Mitte formt einen Beutel, die beiden benachbarten Personen `erbrechen'' symbolisch in den Eimer.
    \item Toaster: Die benachbarten Personen drehen sich zur Person in der Mitte und bilden einen weiten Kreis, indem sie sich an den Händen nehmen. Der `Toast'' in der Mitte springt nach oben und macht `Pling''.
    Waschmaschine: Die beiden äußeren Personen formen mit den Armen einen großen Kreis (vertikal). Die Person in der Mitte `schleudert'' darin ihren Kopf im Kreis.
    \item Fußball: Die Person in der Mitte schießt ein imaginäres Tor, die beiden äußeren Personen jubeln.
    \item  Mixer: Die mittlere Person berührt mit dem Zeigefinger jeweils den Kopf der beiden äußeren Personen. Die beiden äußeren Personen drehen sich jeweils um die eigene Achse.
    \item  Mixer: Die mittlere Person berührt mit dem Zeigefinger jeweils den Kopf der beiden äußeren Personen. Die beiden äußeren Personen drehen sich jeweils um die eigene Achse.
    \item Feuerwehr: Die Person in der Mitte zieht eine imaginäre Leiter hoch, während die beiden benachbarten Personen so tun, als würden sie Wasser aus Feuerwehrschläuchen verspritzen.
    \item Elefant: Die Person in der Mitte formt mit den Händen einen Rüssel vor dem Gesicht. Die beiden äußeren Personen formen jeweils mit ihren Armen die Ohren des Elefanten.
    \item Esel: Alle bleiben stehen und schauen mürrisch. 
\end{itemize}

\subsection{Geheimer Freund}

\begin{itemize}
    \item Dauer: gesamte Akademie
    \item Die Namen aller Gruppenmitglieder werden auf einzelne Zettel notiert, zusammengefaltet und in einen nicht transparenten Behälter gegeben.
    \item Nun ziehen alle einen Zettel und erfahren so ihren stillen Freund, ihre stille Freundin. Es kann nicht getauscht werden, es sei denn, man hat sich selbst gezogen.
    \item Alle werden aufgefordert, ihren Freundinnen und Freunden in der nächsten Zeit kleine, positive Gesten und/oder Aufmerksamkeiten zukommen zu lassen, ohne ihre Identität preiszugeben.
    \item Die Gesten/Aufmerksamkeiten können zum Beispiel sein:
    \item  Ausleihen von Schulmaterialien
    \item aufmunternde Worte
    \item  ein einfaches Stuhlheranziehen, Jacke anreichen oder freundliche Begrüßungen
    \item kleine Zettelchen mit netten Botschaften 
    \item Am Ende des vorher festgelegten Zeitraums (zum Beispiel am Ende einer Projektwoche) wird aufgelöst, wer wen hatte. 
\end{itemize}

\end{document}